\documentclass[12pt]{article}
\usepackage[utf8]{inputenc}
\usepackage[T1]{fontenc}
\usepackage{hyperref}
\title{Kambiw\'a Language Resources for K-12 Curriculum Development}
\author{Compiled by Claude Sonnet 5.5 (claude-sonnet-5-5) for Andr\'as Kornai}
\date{}
\begin{document}
\maketitle

\section{Introduction}
Kambiw\'a (ISO 639-3 xbw; Glottocode kamb1239) is a dormant, unclassified language of Pernambuco, Brazil. Search results describe it as extinct, known only from a few short word lists. The Kambiw\'a people (Inaj\'a, Ibimirim and Floresta, Pernambuco) remember a few words. Sources differ on the dates of the word lists, so those dates are uncertain. Very little exists, and this report does not pad that.

\section{Primary Educational Resources}
I found no primers, textbooks or school programmes for Kambiw\'a in the search results. The language is not known to be taught in schools.

\section{Linguistic and Reference Materials}
\begin{itemize}
\item Wikipedia article on the Kambiw\'a language. It summarises the word lists (Pickering 1961, published in Meader 1978) and Loukotka's classification as unclassifiable: \url{https://en.wikipedia.org/wiki/Kambiw%C3%A1_language}
\item UFRJ thesis record, ``Os \'indios kambiw\'a de Pernambuco: arte e identidade \'etnica'' (Portuguese). It is about art and ethnic identity, not language description: \url{https://bdtd.ibict.br/vufind/Record/UFRJ_7f3f8067158366249d03fc9cabc385aa}
\end{itemize}

\section{Religious and Community Resources}
Joshua Project has a profile of the Kambiw\'a people: \url{https://joshuaproject.net/people_groups/12471}. It is a missions-oriented demographic profile and holds no language materials. I saw no community language organisation in the results.

\section{Digital and Audio Resources}
I found no audio, video, corpus or software resources.

\section{Curriculum Development Context}
The only linguistic data are a couple of dozen words from 20th-century lists. A curriculum can therefore be at most a heritage and cultural-identity programme, not a language-immersion one. Any use of the surviving words should be developed with Kambiw\'a community members. Materials on Kambiw\'a art and territory (for example the UNIVASF paper \url{https://portais.univasf.edu.br/cartes/conartes/anais-do-i-conartes/artigos/kambiwa-em-espalhamentos-provocacoes-sobre-fronteiras-politicas-e-poeticas-de-territorios-indigenas-atraves-do-barro-debora-caroline-viana-almeida}) could serve as cultural context.
\end{document}